% §1 Introduction of the isotropic paper --- new with the split (Q-0307, 2026-10-04; adr/0001),
% assembled from the single article's introduction (paper/01-introduction.tex before the split): its
% second paragraph on the axioms of the line papers, its paragraph "The isotropic corner reduces to
% the line" with the comparison with Schoenberg, its contribution 2, its scope sentence on the line
% theory, and the part of 1.1 that concerns section 3, each moved with the wording the new context
% needs. New: the question of this paper alone, the Matern reading of the threshold (from section
% 3.3), the plan. Every sentence that restates a result was compared with its statement
% (WRITING.md § 5). Nothing here is shared with the blueprint.
%
% 1.1 owns the trust base (article-kit PROCESS.md § 7, WRITING.md § 4): the state of the
% formalization, the four cited facts with the statements that rest on them (the blueprint's
% ledger, blueprint/AXIOMS.md, and the cited-fact sentences of section 3), and the figure.
%
% Q-0310 (the whole-paper review's fixes for this paper; records/affine/
% REVIEW-paper-whole-fixes-paper-iso.md): the strictness of Schoenberg's chain is Gneiting's (1999,
% p. 94, Euclid's hat; READING-gneiting1999radial.md, PRIOR-ART-novelty-claims.md).
%
% Q-0311 (the blind review records/affine/REVIEW-paper-iso.md; its flags answered in
% REVIEW-paper-iso-fixes.md): "hemigroup" defined in the first paragraph, with Hazod-Siebert as
% Paper V cites it (Def. 2.14.16, p. 380; READING-hazod2001stable.md) (flag 5); the comparison with
% the classical axioms stated once, as Paper V states it, against Pauwels et al. (flag 12); the
% kernels from scale zero, never "from the origin" (flag 24); the question framed as a
% classification of kernel families (flag 8); a paragraph on prior work in probability (Sato
% Thm. 16.1, p. 99, by Sato 1991, p. 117; Hazod-Siebert Thm. 2.14.21, p. 383; Sato Thm. 28.4 and
% Lemma 28.5, p. 191, by Zolotarev and Wolfe, Rem. 28.6, p. 193; READING-sato1999levy-prior-art.md)
% and one on the scale-space neighbours (Pauwels et al. p. 695; Duits et al. pp. 267, 277) (flags
% 1, 2, 17); the novelty paragraph and the contributions narrowed to the three new results, and its
% two self-contradictions removed (Schoenberg's Theorem 6 is on the infinitely divisible class, as
% the intersection; the witness family is [B]'s) (flags 6, 7, 10, 22); the reduction sentence kept
% here and in 3.2 only (flag 33); the trust base in plain words, with the cited facts renamed C1-C4
% (flags 3, 25, 28).

%
% Q-0316 (pass 2, the blind register review records/affine/REVIEW-register-paper-iso.md; every flag
% answered in REVIEW-register-paper-iso-fixes.md): the prior-work paragraphs split into short
% sentences and a list (flags 6, 7); (ND) named and the parametrization hedge dropped (flags 9, 10);
% the Matern kernels said to have Matern densities, for gamma > d/2, with no third nu (flags 11,
% 13); Sato's two-sided k named as his (flag 12); the line's cone, the folded profile, the
% subordinated slice and the Cin rays pointed to section 2 where they are first used (flag 1); the
% bridge's hypothesis named (flag 2); the frame and the new results made to agree (flag 3); the
% scale-space neighbours positioned against the isotropy axioms of Koenderink, Florack et al.,
% Florack and Lindeberg, each read at the page (flag 15; READING-koenderink1984structure.md,
% READING-florack1992scale.md, READING-florack1997image.md, READING-lindeberg2011generalized.md);
% Wolfe's theorem as the cited fact C5 and the radially nonincreasing density as a proposition; the
% trust base lists every import and says what the line papers' own trust bases are (flags 4, 14).
%
% Q-0328: prop:origin-lower-bound-d (Q-0272, refereed blind) in "The answer", in contribution 1 (the
% same topic, so still three results) and in 1.1 (it spends C3(iii) only, ledger A3(iii); no new
% cited fact). Each sentence compared with the statement: d >= 2, c finite, almost everywhere.
%
% Q-0344 (pass 6, the author's review of this section; every comment answered in
% records/affine/REVIEW-author-paper-iso.md): rewritten as questions, then answers, in plain
% language. No internal label: the earlier papers are named as a reader would ("the line paper",
% "the cone paper", each with its citation; a numbered place in them is \prog{V}{...} or
% \prog{B}{...}, printed as a located citation, macros.tex). Self-decomposable laws cited where they
% are defined (Sato Def. 15.1, p. 90, and Thm. 15.3, p. 91; Steutel-van Harn Ch. V, (1.1), whose
% stored copy prints no page numbers; READING-steutel2004infinite.md), and the Matern covariance
% (Matern 1960, (2.4.5)-(2.4.7), pp. 17-18; Stein 1999, pp. 31 and 48-49; Guttorp-Gneiting; Porcu et
% al.; READING-matern1960spatial.md, READING-stein1999interpolation.md). The reduction to the line
% explained by the projection of a rotation-invariant vector; the bridge explained as a Brownian
% motion run for a random time, citing the time-causal paper (fagerstrom2026hemigroup, Thm. 7.3) and
% the cone paper's Lem. 4.2 and Prop. 4.5. Zolotarev and Wolfe cited from their own papers, read at
% the page (READING-zolotarev1963structure.md, READING-wolfe1971continuity.md). The three prior-work
% paragraphs open with the question they position the paper against; the detail moved to its
% section: the line threshold with Sato's constant to 2.4 (it was there already), the polar
% criterion's attribution to 3.2, Schoenberg's Theorem 6 to 3.4. Every sentence restating a result
% compared with its statement (WRITING.md section 5).
%
% Q-0377 (the formalization closed: 125 of 125 blueprint nodes \leanok, main at dfafe99): 1.1 is
% rewritten as "What is machine-checked" on Paper V's model, with the cited fact C6 (Schoenberg's
% radial theorem, ledger A10), the headline #print axioms block and the trust-base markers
% `linkage check` reads; the contributions end with one clause on the formalization, and the plan
% counts six cited facts. The fidelity review (Q-0376) is named as pending in 1.1.
%
% Q-0407 (the author's decision on pull request #100): the classification on operators is a
% statement, Corollary cor:isotropic-classification-operators (section 3.2), machine-checked. "The
% answers" and "Scope" name it; 1.1 lists it under "What is proved", in the table's row for
% C3(ii), to which the row for C3(iii) refers, and in the #print axioms block. The abstract is
% unchanged: its sentence "classifies the rotation-invariant kernel families" is the theorem, and
% it is at 249 words.
%
% Q-0408 (pass 6, the author's review of sections 3.3, 3.4 and 4): checked against the new
% openings of 3.3 and 3.4 and the conclusion, so that the three tell one story. Two edits. In "The
% answers", c is introduced first as the power at which the kernel's transform falls off, the
% reading 3.3 opens with (Lemma lem:isotropic-exponent-log(1) with clause (1) of the theorem give
% log mu-hat / log|xi| -> -c), and then as the limit of the polar profile; "the density of its
% small jumps measured on the scale of |x|^{-d}" is gone ("jumps" was never introduced; section 2
% says "displacements"). The plan names section 4 as the conclusion with the open problems (4.1).
% The third answer already asks the question of 3.4 in its terms ("an exponent of the line, read
% as a function of |xi|"), and the abstract is unchanged.
% Second round (the blind register review of the new prose, REVIEW-register-paper-iso-q0408-fixes.md):
% "every such law ... occurs" cites Lemma lem:isotropic-realisation with the theorem, as section
% 3.4 and the conclusion do; "Scope" says what Remark rem:isotropic-line now holds (the comparison
% at d = 1) and that the not-carried list is section 4.1's.
% The same round removed from sections 3.3 and 3.4 every sentence that said which cited facts a
% statement "rests on" (the dispatching session's decision of 2026-10-10: WRITING.md section 3
% gives that to this subsection). Each was first compared with Table tab:trust-base, and 1.1
% gained what the sentences said and the table did not: C3(ii) is read by clause (1) of
% Proposition prop:isotropic-radial-monotone only and clause (2) reads C3(iii) with C2 (as
% the Lean has radial_monotone_family and radial_monotone_law), and the lower
% bound is verified for every law whose transform is the datum's, C3(iii) entering through the
% law of a truncated datum (origin_lower_bound takes a law mu with IsLawOf D mu).

%
% Q-0414 (the repositioning pass, round 1; records/affine/REPOSITIONING-paper-iso.md has every
% claim of novelty or attribution made here with the statement that carries it, the source, its
% anchor and the reading note). The section is rewritten against the prior art read after the
% paper was written (records/affine/PRIOR-ART-origin-threshold-reading.md): the questions carry
% their stakes for a measurement; "transform", "exponent", "Gaussian part" and "Levy measure" are
% said at first use; the second answer is in frequency-side words (c is the power at which the
% transform falls off, Corollary cor:origin-profile-exponent(2)); the dimension enters twice,
% through the class A_d and through the threshold d; one name for the Matern kernels, introduced
% once with their identification as the symmetric variance-gamma laws; the prior work on the
% centre of the kernel (the line: Zolotarev, Wolfe, Sato for continuity, Sato-Yamazato for
% boundedness and for the rate, Hartman-Wintner; R^d: Watanabe, Knopova-Schilling,
% Bogdan-Grzywny-Ryznar, Grzywny-Ryznar-Trojan, Cho-Kang-Kim); the cone paper positioned in the
% paragraph on the dimension filtration (cone paper, Prop. 4.9, Lem. 4.12(3), Rem. 4.10 and
% Figure 4, read at the page images of pp. 34, 38 and 39 for this item); the contributions as
% the queue item states them. 1.1 keeps the statement of the trust base; its inventory is
% Appendix B (appendix-verification.tex).
% Anchors: grzywny2019asymptotic by statement number of the published text (advance-access
% pagination 1-66; the journal pages 7193-7258 are not verified). (cho2021estimates and
% bogdan2014density, named above, are no longer cited: the rewrite of Q-0428 removed them.)
% Q-0424 (records/affine/reviews/VERIFY-round1-claims-and-trust-base.md, item 9): the paper
% cites grzywny2019asymptotic by statement number only and never by page, so nothing waits on its
% journal folios. One number collides with the first arXiv version, Corollary 5.8 (there Theorem
% 5.8 is another statement and the corollary is absent; READING-grzywny2019asymptotic.md,
% section 7), and Table tab:in-print says "of the published article" at that citation.
%
% Q-0424 (batch R1 of the response to the external review of round 1; records/affine/reviews/
% RESPONSE-round1-R1.md has every finding with what was done). In this file: question 2 says
% what an unbounded kernel does to the record of a point source, in the invariant form (no
% version of the density bounded near the centre), where it said "no finite value to sample"; and
% section 1.1 (its opening, the departures of C5 and C6, the paragraph on C5, the caption of
% Table tab:trust-base, the closed form of A_beta under "Outside the formal verification"). The abstract, "The
% answers", the prior-work paragraphs and the contributions are not touched: they are the
% framing batch's (R4), written once after the new statements are in.
%
% Q-0428 (the rewrite after the external review, round 1; records/affine/
% REPOSITIONING-paper-iso-round2.md has every attribution made here with its anchor and note, and
% every sentence on what the paper contributes with the statement that carries it). The author's
% direction of 2026-10-12 (records/affine/REVIEW-author-paper-iso.md, last section): the paper is
% a paper in scale-space theory; no apparatus of novelty; attribution exact and made once. So:
% the opening says why the semigroup axiom is relaxed and what that admits (the presentation
% review, item 1); mu_{0,t} and "from scale zero" are defined at first use, with "centre" for the
% point x = 0 (item 5); the answers carry the eight statements of Q-0425 (the rate through the
% polar profile, the dictionary and its consequences, the limit of the critical values) and say
% each result once; the Cin rays are introduced as elementary exponents, with the sentence on
% their superpositions and Jurek and Vervaat's Theorem 7.2 (item 14;
% READING-jurek1983integral-driving-process.md); the bridge is a Gaussian variance mixture before
% it is a process (item 15); "Scope" stands with the answers and names what is not treated (item
% 2). The three paragraphs "Prior work" and the list of contributions in two tiers are replaced
% by "Relation to scale-space theory", "Relation to probability" and "What the paper
% contributes". The anchors of the probability literature are in Table tab:in-print, once.
% 1.1, as first written in this item: two lists, machine-checked and not, with three statements
% written as machine-checked ahead of the Lean of the parallel item Q-0427. That Lean is merged
% since, and the fixes round (below) rewrote the subsection.
% cygan2017asymptotic and leonenko2013tauberian are cited in the numbering of their arXiv
% versions and bingham1997abel in that of the Hokkaido preprint: these are the copies held and
% the ones the paper cites (the author, 2026-10-12), said once, in the caption of Table
% tab:in-print. Their published numbering was not compared.
%
% Q-0428, the fixes round (records/affine/REPOSITIONING-paper-iso-round2-fixes.md has every flag
% of the register review of the front and every row of the scope audit of 2026-10-12 with what
% was done). The order of section 1: the cascade, what the semigroup axiom adds and what
% dropping it gives up and gains; the paper's question, with the axioms on kernels in one
% sentence, on the first page; what is known on the line, with the definitions the answers need;
% what the weaker axiom admits; the questions; the answers, each opening with its claim and held
% to the level of the questions (the two forms of the rate with their constants, the formula
% and the conditions of the dictionary, the sentence on the Urbanik representation and the limit
% of the critical values are in the leads of sections 3.3 and 3.4 and in the sources); scope;
% the two relations; the contributions as three items with their sections; the plan.
% Each sentence on the scale-space literature is held to its reading note (the scope audit,
% section 5): scale invariance is named with rotation invariance; Koenderink's route is a
% diffusion equation, not a semigroup; "stable" is this paper's word for the kernels of Pauwels
% et al. and of Florack et al.; Florack et al. are cited for the words of p. 380, whose formulas
% the note did not read on the page image; Florack's book with its p. 57 for "positive Schwartz
% filters". No sentence places the Matern kernels in the scale-space literature: no reading note
% supports one (the register review's flag 12 is a lead for the librarian).
% Section 1.1: three opening sentences (what Lean checked, what it did not, who compared the
% formal statements with the printed ones), the two lists with each print-only statement named
% in words, what is assumed, then the caveats, short, with the accounts statement by statement
% in Appendix B. The second list had four statements when the Lean of Q-0429 was merged, and
% has five since Q-0448 (Proposition prop:step-exponent-dimensions).

\section{Introduction}\label{sec:introduction}


% Q-0448 (the response to the second external round; records/affine/reviews/
% VERIFY-round2-scale-space-literature.md, sections 2, 3 and 10.2). With rotation and scale
% invariance the semigroup axioms leave the transforms e^{-t|xi|^p}, and positivity leaves
% 0 < p <= 2, the Gaussian included; a further requirement singles the Gaussian out (the order
% was the other way round). The two-scale cascade has its precedent, the causality axiom of
% Alvarez, Guichard, Lions and Morel (p. 203, at the page image; READING-alvarez1993axioms-
% architecture.md). "What is given up with the semigroup is an evolution equation" is gone:
% they obtain one without a semigroup, from locality, which is what is not asked here.
% Q-0448, the blind draft review after the split (records/affine/
% REVIEW-draft-paper-iso-q0448.md, with its dispositions beside it), group 2: one telling in
% section 1. The cascade and the aperture are defined in the first paragraph (flag 19); "stable"
% has its range where the kernels are first written (flag 13; Sato's Theorem 14.14 is for
% 0 < alpha <= 2, the Gaussian included, and so is clause (3) of Proposition
% prop:isotropic-bridge); the first paragraph points to "Relation to scale-space theory" for
% the causality axiom of Alvarez et al. and for what singles the Gaussian out, each of which is
% said there, once (flags 11, 12); that no generator is claimed is a limitation, in section 4
% (flag 7). The second paragraph states the axioms and leaves the questions to their list
% (flag 20).
A scale-space representation measures a signal through an aperture whose size is the scale. A
measurement is convolution with a probability kernel, and the record at a coarser scale is
computed from the record at a finer one by a further convolution, a \emph{stage}. The stages
together are the \emph{cascade}, and the \emph{aperture} at scale $t$ is the kernel that takes
the signal itself to its record at scale $t$.
Scale-space axioms usually require the cascade to be a semigroup in one scale parameter, as those of
the extended class of Pauwels, Van Gool, Fiddelaers and Moons do \sscite{pauwels1995extended}. A
semigroup requires more than a cascade needs: every stage has to be itself one of the apertures, so
that one parameter indexes the stages and the apertures alike. With rotation and scale
invariance the semigroup axioms of the literature leave the probability kernels whose Fourier
transform at the frequency $\xi \in \RR^d$ is $e^{-t|\xi|^p}$, for one $p$ with $0 < p \le 2$.
This paper calls them the \emph{stable kernels} of index $p$. The Gaussian is the stable kernel
of index $2$. The axioms studied here require only that the stages compose. With
$\mu_{s,t}$ the kernel of the stage from scale $s$ to scale $t$, this is
$\mu_{r,t} = \mu_{s,t} * \mu_{r,s}$ for $r \le s \le t$, and how a stage depends on its two
scales is left free. Such a family is a \emph{convolution hemigroup}
(\sscite{hazod2001stable}, Def.~2.14.16, p.~380). What is gained is a larger class of
apertures. The two-scale cascade has a precedent in scale-space theory, and ``Relation to
scale-space theory'' below sets the present axioms beside it and beside the semigroup axioms.

The axioms are those of the author's paper on the line \sscite{fagerstrom2026spatial}, the
\emph{line paper} below, with rotation invariance in place of the reflection symmetry
(Definition~\ref{def:isotropic-operator-family}): an aperture has a size, and no shape and no
preferred direction. Read on the kernels, they say that the stages
compose, that they depend continuously on the pair of scales, that each is rotation invariant,
and that a dilation of the signal only relabels the scales
(Definition~\ref{def:isotropic-kernel-family}). A family of kernels with these properties is an
\emph{isotropic cascade family}, an \emph{isotropic family} for short. The axioms read the same
in every dimension. In imaging $d$ is $2$ or $3$.

On the line the families are classified in the line paper
(Theorem~\ref{thm:line-classification}), and a second paper \sscite{fagerstrom2026cone}, the
\emph{cone paper}, describes the set of them. What is used from the two is restated in
\S\ref{sec:line}. A kernel is a probability measure, and that classification, like everything
after it in this paper, rests on the theory of such measures in probability. So from here on
the paper uses the vocabulary of that theory: a kernel is a \emph{law}, and its Fourier
transform, the transfer function of the smoothing, is the characteristic function of the law.
Stage, scale and cascade keep their meaning, and ``kernel'' stays in use for a law in its role
as an aperture. A frequency is written $\omega$ on the line and $\xi$ in $\RR^d$, and
$\widehat\mu(\omega) = \int e^{i\omega x}\,\mu(dx)$ is the Fourier transform of a kernel $\mu$
on the line, its \emph{transform} for short. The axiom on dilations
makes the apertures, the kernels $\mu_{0,t}$ \emph{from scale zero}, the dilations of one law:
up to a relabelling of the scales, $\mu_{0,t}$ has the
transform $e^{-F(t|\omega|)}$ for one function $F$, the \emph{exponent} of the family. The
cascade makes that law \emph{self-decomposable}: for every $0 < \lambda < 1$ it is the
convolution of its own image under $x \mapsto \lambda x$ with another law
(\sscite{sato1999levy}, Def.~15.1, p.~90; \sscite{steutel2004infinite}, Ch.~V, (1.1)). In a
cascade this is the relation between two scales. The kernel at the finer scale is a contracted
copy of the kernel at the coarser one, and the other law is the stage between them. Every
symmetric self-decomposable law other than the point mass occurs. Such a law is the convolution
of a centred Gaussian, which may be absent, with the law of a sum of independent displacements.
A law of the second kind alone has \emph{no Gaussian part}. Its \emph{\Levy{} measure} gives the
expected number of displacements of each size. On the line the density of that measure at the
size $x$, the two directions counted together and multiplied by $x$, is the \emph{folded
profile} of the law. In $\RR^d$ the same quantity is
taken in polar coordinates, the \emph{polar profile} of \S\ref{sec:isotropic-reduction}. The
Gaussian part and the profile determine the exponent, and the exponent determines them
(\S\ref{sec:line-axioms}; Proposition~\ref{prop:isotropic-polar-reading}(1)).

% Q-0448, the blind draft review (flags 12 and 13): what the semigroup axioms arrive at is said
% in the first paragraph and, with its sources, under "Relation to scale-space theory"; here
% the stable kernels are named as laws. The Bessel scale space is described there, once.
This is what the weaker axiom admits. In probability the stable kernels are the
rotation-invariant stable laws (\sscite{sato1999levy}, Thm.~14.14, p.~86). Among the kernels
added are the laws on the line with transform
$(1 + \omega^2/2)^{-\gamma}$, $\gamma > 0$, whose folded profile is $2\gamma e^{-\sqrt2\,x}$.
This paper calls these laws and their dilations the \emph{\Matern{} kernels} of index $\gamma$.
Scale-space theory knows them as the kernels of the Bessel scale space (``Relation to
scale-space theory'' below). On the line the kernels of integer index are those computed by a
cascade of identical first-order recursive filters, run forward and backward along the signal
(\prog{B}{Prop.~7.2 and Cor.~5.5}). Read as functions of the length $|\xi|$ of a frequency
$\xi \in \RR^d$, the transforms of the stable and the \Matern{} laws are the
transforms of self-decomposable laws on $\RR^d$, for every $d$
(Propositions~\ref{prop:line-filtration}(2) and~\ref{prop:isotropic-bridge}). For
$\gamma > d/2$ the density of the \Matern{} kernel on $\RR^d$ is a \Matern{} covariance
function, of smoothness $\gamma - d/2$. Spatial statistics uses that family to model random
fields whose smoothness is a parameter, and \Matern{} obtained it as a gamma mixture of
Gaussians (\sscite{matern1960spatial}, (2.4.5)--(2.4.7), pp.~17--18;
\sscite{stein1999interpolation}, \S2.7, p.~31, and \S2.10, pp.~48--49;
\sscite{guttorp2006studies}; \sscite{porcu2024matern}). In probability the \Matern{} kernels are
the symmetric variance-gamma laws \sscite{fischer2025variance}, and in potential theory their
densities are dilations of the Bessel kernel \sscite{aronszajn1961theory}
(Example~\ref{ex:matern} has the formulas).

\paragraph{The questions.} Three questions arise.
\begin{enumerate}
\item Which kernel families do the axioms allow in $\RR^d$? Their kernels from scale zero are
the apertures that a scale space with these axioms can use in the plane and in space.
\item What does the dimension change in a kernel? Chiefly, is the kernel bounded at its
\emph{centre}, the point $x = 0$, or does it record an ideal point source as a spike?
\item Does the exponent of a family of the line, read as a function of $|\xi|$, give a family in
the plane and in space? This is the rotation-invariant extension of the family, not the
separable one. The Gaussian and the $\alpha$-scale spaces of Duits et al.\
\sscite{duits2004axioms} pass. Is that general,
or does the set of allowed exponents shrink as $d$ grows?
\end{enumerate}

% Q-0448 (the author's decision of 2026-10-13: what is said twice goes; exposition item 4 of
% records/affine/reviews/VERIFY-round2-length-and-exposition.md): "The answers" and "What the
% paper contributes" are one passage. It opens with the main result, as the author's section
% review of 2026-10-12 asked of the contributions, and the two descriptions follow as the second
% and the third answer.
\paragraph{The answers.} The main result answers the first question. A probability law on
$\RR^d$ is a kernel from scale zero to a positive scale of an isotropic family, an aperture of a
scale space with these axioms, exactly when it is
rotation invariant, self-decomposable and not the point mass at the origin. In every isotropic
family the kernels from scale zero are, up to a relabelling of the scales, the dilations of one
such law, and every such law gives an isotropic family, for every relabelling of the scales
(Corollary~\ref{cor:isotropic-scale-zero-kernels}, with
Theorem~\ref{thm:isotropic-classification} and Lemma~\ref{lem:isotropic-realisation}; on the
operators the theorem is Corollary~\ref{cor:isotropic-classification-operators}). The reason is
a reduction to the line.
A rotation-invariant kernel is determined by its projection on a line through the origin, and
the projections of an isotropic family form a family on the line
(Lemmas~\ref{lem:radial-marginal} and~\ref{lem:isotropic-marginal-family}). Each aperture has
a radially nonincreasing density, so it weights a nearer point at least as much as a farther one
(Proposition~\ref{prop:isotropic-radial-monotone}).

Second, the dimension is a threshold for the kernel at its centre. A kernel with a Gaussian part
is bounded and smooth in every dimension: its transform falls off faster than every power
(\S\ref{sec:isotropic-origin}). For a kernel from scale zero with no Gaussian part,
let $c$ be the power at which its transform falls off at high frequencies:
$-\log\widehat\mu(\xi)/\log|\xi| \to c$ as $|\xi| \to \infty$, with $0 < c \le \infty$
(Corollary~\ref{cor:origin-profile-exponent}). The kernel has a density that is bounded near its
centre exactly when $c > d$, and for $c > d + N$ the density is sure to have $N$ continuous
derivatives (Theorem~\ref{thm:origin-boundedness-d}). The number $c$ belongs to the exponent. It
is the same in every dimension that admits the exponent, and only the threshold moves. So one
exponent can give a kernel that is bounded on the line and has a spike in the plane, and the
\Matern{} kernels, for which $c = 2\gamma$, do this for $\tfrac12 < \gamma \le 1$. When $c < d$
the density grows at the centre like $|x|^{c-d}$ times a \emph{slowly varying} factor, one that
changes more slowly than every power, as a constant or a logarithm does, and when $c = d$ it
grows at least like $\log(1/|x|)$. The factor and the constant are computed from the \Levy{}
measure (Proposition~\ref{prop:origin-rate-profile-d}). Its small displacements decide whether
the kernel is bounded at its centre and how many derivatives it is sure to have, and for a
spike its power and its slowly varying factor. The large ones enter the constant of the spike.

% Q-0448, the blind draft review (flag 24): the stages have a paragraph of their own, with both
% cases of Proposition prop:isotropic-stages (a point mass exactly for no Gaussian part and a
% finite c; a density otherwise). The sentence on what a designer specifies is said of the Levy
% measure, which is what the statements take as given.
A stage of the cascade need not be an aperture. When the kernels have no Gaussian
part and $c$ is finite, the stage from scale $s$ to scale $t > s > 0$ passes a fixed fraction
of the finer record through unchanged and blurs the rest. The fraction is $(s/t)^c$ when the
scales are labelled so that the kernel at scale $t$ is the dilation by $t$. With a Gaussian
part, or with $c = \infty$, as for the stable kernels, a stage is a blur by a density
(Proposition~\ref{prop:isotropic-stages}).

% Q-0448, the blind draft review (flags 2, 21, 22, 23): the third answer at the level of its
% question, in four points: strict at every dimension; what remains in all of them, said as
% Proposition prop:line-filtration(3) says it (the mixtures whose mixing law is a kernel of a
% time-causal family, not every Gaussian variance mixture: every member of the separating
% family is one and is admitted in finitely many dimensions); the simplest exponents, in the
% two sentences of the author's section review of 2026-10-12, unchanged; the test, said and not
% announced, and the family. The detail, with the time-causal literature (now cited in section
% 2.3, where the time-causal theory is introduced) and Proposition
% prop:step-exponent-dimensions, is in sections 2.3 and 3.4.
Third, each dimension admits strictly fewer exponents than the one before. The sets
$\mathcal{A}_d$ of the exponents admitted in dimension $d$ form a strictly decreasing sequence,
the \emph{dimension filtration} (Propositions~\ref{prop:line-filtration}
and~\ref{prop:profile-one-dimension-down}). The exponents that remain in every dimension are
those of the Gaussians whose variance is drawn at random from a kernel of a time-causal scale
space with the same axioms, one that measures a signal in time from its past alone. These
Gaussian variance mixtures include the stable and the \Matern{} kernels
(\S\ref{sec:line-bridge};
Propositions~\ref{prop:line-filtration}(3) and~\ref{prop:isotropic-bridge}). Every exponent of
the line is a
superposition of a Gaussian term and of elementary exponents, each with a \Levy{} measure whose
density is proportional to $1/|x|$ up to a cut-off and zero beyond it
(Proposition~\ref{prop:line-extreme-rays}). The exponents admitted in the plane are still such
superpositions, the \Matern{} and the stable ones among them. What is lost first is each
elementary exponent taken alone: none of them is admitted in the plane
(Proposition~\ref{prop:profile-one-dimension-down}). For the next dimension there is a test.
The exponent of a kernel with no Gaussian part and a finite $c$ determines a function of the
radius $|\xi|$ that is the Fourier transform of a probability law on $\RR^d$, and the exponent
is admitted in $\RR^{d+1}$ exactly when that function is still the transform of a law there
(Proposition~\ref{prop:radial-dictionary}). And one explicit family of exponents is followed
through the dimensions: the dimensions that admit a member are read off its parameter
(Theorem~\ref{thm:separating-family-critical-values}).

% Q-0448, the blind draft review (flag 7): the scope says what the paper is about. Its
% limitations (the kernels from scale zero, the smoothing part of a stage, locality, the
% extrema, the generator) are listed once, in section 4.
\paragraph{Scope.} The axioms and the classification are stated on operators and on their kernels,
and the two forms describe the same families (Lemma~\ref{lem:operators-to-kernels}). The other
statements are on kernels. The paper is theory: it has no
discretization, no implementation and no experiment, and \S\ref{sec:conclusion} lists its
limitations. Apertures with a shape are the subject of
a companion paper on affine covariance, not yet released.
% TODO(module affine): cite the affine paper of this repository (paper/) when released.

% Each sentence against its reading note (records/affine/READING-koenderink1984structure.md,
% -florack1992scale.md, -florack1997image.md, -lindeberg2011generalized.md,
% -pauwels1995extended.md, -duits2004axioms.md; the scope audit of 2026-10-12, section 5).
% Lindeberg is cited by section; the held copy is the preprint, whose pages 5 and 10 these are.
% Florack et al.: (6) and (19), (23)-(25) were read at the page images of pp. 379 and 380
% (READING-florack1992scale-p380.md, Q-0430): G(omega; sigma) = exp(alpha sigma^p omega^p),
% alpha an arbitrary negative constant, "only positive p-values will be of interest", omega the
% magnitude of the frequency; no upper bound on p is stated there and positivity of the kernel
% is not among their conditions, so "when 0 < p <= 2" in the first sentence is this paper's
% (the range in which the filter is the transform of a law; Duits et al. below have it for
% their generators). Pauwels et al., p. 694: recursivity and scale invariance give the transforms
% e^{-a|omega|^alpha t}, alpha > 0; the word "stable" is not theirs.

% Q-0448 (the response to the second external round; VERIFY-round2-scale-space-literature.md,
% sections 1 to 3 and 10.3, with the page of every anchor and the reading note that read it:
% READING-alvarez1993axioms-architecture.md, pp. 203, 204, 206, 211, 214, 226;
% READING-felsberg2004monogenic.md, pp. 7-8; READING-weickert1999linear.md, pp. 239-240;
% READING-burgeth2005bessel.md, pp. 86, 89, 90; READING-burgeth2005relativistic.md, pp. 3, 7).
% Four short paragraphs under one heading: the two-scale cascade and the scale invariance with
% a free relabelling are axioms of Alvarez et al., who ask locality and regularity besides; the
% semigroups with scale invariance, and what leaves the Gaussian alone; the semigroups without
% scale invariance, the Bessel scale space among them, whose kernels are the Matern kernels
% and whose dilations at a fixed index have no composition law there; what the present axioms
% give. Not written, each for the reason of section 10.7 of the record: that the family of
% Pauwels et al. was derived by Iijima; that the shadow is Matheron's descente; that the
% recursive Gaussians are Matern kernels; any constant of p. 87 of Burgeth et al.
\paragraph{Relation to scale-space theory.} The two-scale cascade is the causality axiom of
Alvarez, Guichard, Lions and Morel, who state it as the weaker form of the semigroup, and their
scale invariance, like axiom (A8), lets a dilation relabel the scales in a way that is not
prescribed (\sscite{alvarez1993axioms}, pp.~203, 206 and~226). They require in addition that the
stages be regular and local. For linear operators that commute with the isometries the heat
equation follows, up to a rescaling of the scale (Thm.~3, p.~214). The present axioms keep their
comparison principle as positivity and do not require locality.

% Q-0448, the blind draft review (flags 12, 13, 14, 38). The range of "stable" is fixed in the
% first paragraph, and p = 2 alpha is said for Duits et al. Each route to the Gaussian has a
% sentence that names its assumption, from the reading note of its source: Iijima as reported
% by Weickert et al. (the transforms exp(-xi^(2m)), m = 1, 2, ..., on the line, Lemma 4 and
% Theorem 1, p. 240; positivity gives m = 1) with the regularity that Felsberg and Sommer name
% (p. 8); Koenderink, p. 365 (spurious resolution prohibited, then homogeneity and isotropy);
% Florack et al., (23)-(25), p. 380 (separability fixes p = 2); Florack's book, Assumption 3.1
% and Observation 3.1, p. 57, and Result 3.2, p. 61; Lindeberg, sections 3.1 and 4.1.1. The
% paragraph that recapped the three answers is gone.
The axiom systems that keep a semigroup with scale invariance arrive at the transforms
$e^{-a|\xi|^p}$: Pauwels et al.\ from recursivity and scale invariance
(\sscite{pauwels1995extended}, p.~694 and \S III, p.~695), Duits et al.\ as the semigroups
generated by $-(-\Delta)^\alpha$, positive for $\alpha \in (0,1]$, where $p = 2\alpha$
(\sscite{duits2004axioms}, pp.~267 and~277), with the Poisson scale space of Felsberg and
Sommer at $\alpha = \tfrac12$ (\sscite{felsberg2004monogenic}, Thm.~1, p.~7). A requirement of
another kind leaves the Gaussian alone, and each derivation has its own. Iijima's axioms, the
first, give on the line the transforms $\exp(-\xi^{2m})$, $m = 1, 2, \dots$, of which
positivity leaves the Gaussian (\sscite{weickert1999linear}, pp.~239--240). His derivation
supposes a transform that is differentiable at the origin, which that of the Poisson kernel is
not (\sscite{felsberg2004monogenic}, p.~8). Koenderink prohibits spurious resolution, his
causality, and with homogeneity and isotropy obtains the diffusion equation
(\sscite{koenderink1984structure}, p.~365). Florack et al.\ require separability
(\sscite{florack1992scale}, (23)--(25), p.~380). Florack's book requires filters that are positive
Schwartz functions and form a convolution algebra (\sscite{florack1997image}, p.~57 and
Result~3.2, p.~61). Lindeberg requires that local extrema not be enhanced
(\sscite{lindeberg2011generalized}, \S3.1 and \S4.1.1).

Without scale invariance a semigroup of rotation-invariant kernels need not be of this kind. The
Bessel scale space of Burgeth, Didas and Weickert is the semigroup in $t \ge 0$ of the kernels
with transform $(1 + |\xi|^2)^{-t/2}$, which at $t = 2\gamma$ are dilations of the \Matern{}
kernels of index $\gamma$ (\sscite{burgeth2005bessel}, pp.~86 and~89). It runs along the index
and is not scale invariant. Its authors also dilate the kernels at a fixed index, with no rule
for composing the dilations (p.~90). The hemigroup is that rule
(Proposition~\ref{prop:isotropic-bridge}). The relativistic scale spaces of the same authors
are a second such semigroup \sscite{burgeth2005relativistic}.

% The classification as a statement about laws, and the statements on the centre and on the
% filtration: every anchor is in the table of the sources (sources.tex), once. This paragraph
% says in words what is in print and what the proofs use (the scope audit, row 5: the proofs use
% the seven cited facts, and "uses" was a claim about them that the trust base denies).
% Q-0448, the blind draft review (flags 15 and 40): the two theorems are named by their numbers
% at the claim (their pages and what each states are in Table tab:in-print, which the paragraph
% points to); the clause on the representation of a self-decomposable law by an integral of a
% Levy process, which named no statement of the paper, is gone (the table has it at the
% dictionary). The count of the cited facts is said in the abstract, in 1.1 and in 2.6, and
% no longer here and in the plan.
\paragraph{Relation to probability.} The laws that the axioms select are a known class. Read on
laws, the classification is in print, for the self-similar processes with independent
increments (\sscite{sato1999levy}, Thm.~16.1) and for hemigroups on groups
(\sscite{hazod2001stable}, Thm.~2.14.21). In both the action of a
dilation on the parameter is prescribed. Here it is a relabelling of the scales that the axioms
leave free. The statements on the kernels have counterparts in print as well, for a case or in
another form: on the line, results of Zolotarev, of Wolfe and of Sato and Yamazato
(\S\ref{sec:line-origin}); in $\RR^d$, Tauberian theorems, which pass from the decay of a
transform to the growth of a density. Table~\ref{tab:in-print} in
\S\ref{sec:isotropic-sources} says, statement by statement, where each is in print and for
which case. The proofs given here are written out from the cited facts of
\S\ref{sec:cited-facts} and the statements restated in \S\ref{sec:line}.

\paragraph{Plan of the paper.} \S\ref{sec:line} restates what is used from the line paper and the
cone paper, and the theorems of the literature that the proofs cite, the
\emph{cited facts}.
\S\ref{sec:isotropic} is the isotropic theory: the axioms, then the three questions in their
order.
\S\ref{sec:conclusion} has the limitations and what is left open.
Appendix~\ref{app:isotropic-proofs} has a table of the notation (Table~\ref{tab:notation}), the
sources (\S\ref{sec:isotropic-sources}), the proofs of the classification, the statements that
\S\ref{sec:isotropic} names and does not state, and the proofs of the statements that are
proved in print only. Appendix~\ref{app:verification} says where the formal development is. The
other proofs that Lean checked, the cited facts in full and the list of the formal declarations are
in \suppl. \supplnumbers{} What is machine-checked, and what is proved in print only, is said
next.

\subsection{What is machine-checked}\label{sec:trust-base}
% trust base: begin
% This subsection is the STATEMENT of the trust base (article-kit WRITING.md section 4): what is
% machine-checked, what is assumed, what is neither. Its inventory (tool versions, commands, the
% Lean names of the cited facts, the table of the cited facts per statement, the stamped #print
% axioms block, the account of the statements of section 2 one by one) is Appendix B,
% appendix-verification.tex, under a second pair of "trust base" markers. `linkage check` reads
% every marked region of every .tex file of the paper directory (LINKAGE.md rule 13), each on
% its own. It reads a row of a table, an \item or a sentence that \ref's a statement and says
% "machine-checked" as a claim that the node is \leanok, and "not machine-checked" as a claim
% that it is not. So the items of the two lists are checked one by one, and no other sentence of
% the subsection gives a statement a status.
% C1-C4 are the ledger's A1-A4, C5 is A8, C6 is A10, C7 is A11 (blueprint/trust-boundary.txt).
%
% Q-0428, the fixes round (the register review of the front, structural change 3; the author's
% answer of 2026-10-11, records/affine/REVIEW-author-paper-iso.md: "I have not reviewed the Lean
% statements myself"). The opening is three sentences: what Lean checked, what it did not, who
% compared the formal statements with the printed ones. The print-only statements are named in
% words beside their numbers. The state the prose is written for: five statements proved in
% print and not machine-checked (lem:exponent-profile-constant, prop:radial-dictionary,
% cor:dictionary-filtration, prop:separating-family-mixtures, and since Q-0448
% prop:step-exponent-dimensions), every other statement machine-checked.
%
% Q-0431 (the author's section review of 2026-10-12, verbatim in
% records/affine/REVIEW-author-paper-iso.md: "we focus on what is not what isn't, and this is
% discussed in the AI statement"; and, of the sentences on Proposition prop:isotropic-bridge,
% "It gives the impression that we have just invented an hypothesis out of thin air. When we in
% fact have cited a published result"). The subsection says what is the case. The opening: what
% Lean checked; where the seven theorems are stated; who compared the formal statements with
% the printed ones, with a pointer to the statement on the use of AI; how the rest is
% established (in print, in the text, by scripts). "It did not check ..." and "no human, the
% author included, has read the formal statements" are gone, and so is "No human has read a
% formal statement" at the end of the paragraph on the review. The two lists stay as they are:
% they say what is, and `linkage check` reads them. After the lists, Proposition
% prop:isotropic-bridge is said to be stated for the kernels of a time-causal family, whose
% existence for the three named delays is two cited results (the time-causal paper's
% classification, Halgreen's theorem), outside the formal development: the formal statement
% takes the family as given (IsCausalKernels, and hT for the inverse-gamma delay). "Conditional"
% and "assumes" are gone. The count of seven assumed theorems is about the formal statements
% and stays true. Other sentences turned from what was not done to what is: the statements of
% section 2 in their own papers; the three things "neither machine-checked nor cited" (they are
% proved in print or computed); "No statement uses a computed value"; "no profile with such a
% factor is constructed in Lean" (the example of section 3.3 is computed in the text).
% Q-0448 (the split, specification 3(b)): the subsection is one paragraph and the two lists.
% Its other paragraphs (the statements of section 2, the form of five formal statements, what
% is outside the formal verification, the review of the formal statements) are section S1 of the
% supplement, supp-trust.tex, cut whole (records/affine/split_q0448_roots.py).
% Q-0448, the blind draft review of the split (records/affine/REVIEW-draft-paper-iso-q0448.md,
% flags 1, 3 and 40). The opening names the three parts whose statements have formal
% counterparts, as it did before the split ("of this paper and of its supplement" took in the
% restatements of section 2 and the cited facts of section S2). The three caveats that the
% split had left in section S1 alone stand after the lists (WRITING.md section 4: nothing
% honesty-critical leaves 1.1): the statements of section 2, the time-causal families of
% Proposition prop:isotropic-bridge in the wording of Q-0431, and what is computed. Lemma
% lem:filtration-top is named with what it states. The closing sentence, which repeated the
% plan, is gone.
The statements of \S\ref{sec:isotropic}, of Appendix~\ref{app:isotropic-proofs} and of
\S\ref{supp:proofs} have formal counterparts in the Lean~4 proof
assistant, on its mathematical library Mathlib, except those of the second list below. Lean
checked that each formal statement follows from the formal definitions, its own logical
foundations and seven theorems of the literature, the cited facts \ledger{A1}, \ledger{A2},
\ledger{A3}, \ledger{A4}, \ledger{A8}, \ledger{A10} and \ledger{A11}, which are its only
assumptions. \S\ref{sec:cited-facts} names the seven, each with the page of its source, and
says with which statement each enters: Theorems~\ref{thm:isotropic-classification}
and~\ref{thm:origin-boundedness-d}, Propositions~\ref{prop:isotropic-polar-reading},
\ref{prop:isotropic-radial-monotone} and~\ref{prop:origin-rate-profile-d}, and
Lemma~\ref{lem:filtration-top} of \suppl.
% Q-0449 (the author's final reading, 2026-10-13): the sentence that said with which statement
% each cited fact enters (C3 with the classification, ...) was a forward reference to section
% 2.6; it is there now, after the list. The statements are still named here, without the
% C-numbers: LINKAGE.md rule 13 compares the cited facts this region names with those that the
% statements it \ref's read, and two of them (A8, A11) are read by no other statement named in
% this section.
Table~\ref{tab:trust-base} in
Appendix~\ref{app:verification} has the cited facts of every statement. The statements of the second list are proved in
\S\ref{app:print-proofs}, and their proofs were refereed by AI agents that saw the statements
and their proofs and no other part of the work. The formal statements were compared with the
printed ones by AI agents. \S\ref{supp:trust} has the account of that review, of the form of
the formal statements and of what the text computes, and the statement on the use of AI at the
end of the paper says how the work was done.
\begin{itemize}
% `linkage check` reads each \item as one row and compares its status with the \leanok tag of
% every node it names.
\item \emph{Machine-checked}: every definition, lemma, proposition, theorem and corollary of
these three parts that the second list does not name.
\item \emph{Not machine-checked}, proved in print:
Proposition~\ref{prop:radial-dictionary}, the dictionary between the exponents of logarithmic
growth and the radial characteristic functions;
Corollary~\ref{cor:dictionary-filtration}, the two readings of the test in that dictionary;
Proposition~\ref{prop:separating-family-mixtures}, the separating family as laws;
Lemma~\ref{lem:exponent-profile-constant}, the closed form of the constant $\varkappa_d$;
and Proposition~\ref{prop:step-exponent-dimensions}, the step exponent one and two dimensions up.
% Q-0448, the joining step: Proposition prop:isotropic-stages stood in this list until the Lean
% of queue item Q-0447 was merged. Its three blueprint nodes are \leanok, it falls under the
% first item, and its proof is in the supplement (supp-proofs.tex inputs proof-stages.tex).
% Five statements are in this list.
% Q-0428, at the merge of the Lean of Q-0429 (its second part): the sentence "At this revision
% also: ..." that stood here, with Lemma lem:shadow-polar-profile, Proposition
% prop:profile-one-dimension-down and Theorem thm:separating-family-critical-values, is gone.
% The three are machine-checked and fall under the first item, which names no statement: it is
% the complement of this one. Their items, with their declarations, are in the list of Appendix
% B, where `linkage check` reads their status.
\end{itemize}
The statements of \S\ref{sec:line} are restated from the line paper and the cone paper. The
formal development proves what it needs of them from the seven cited facts, and in their own
papers part of them is proved in print. Proposition~\ref{prop:isotropic-bridge} is stated for
the kernels of a time-causal family. That each of its three named delays has such a family is
the classification of the time-causal paper, with Halgreen's theorem for the inverse-gamma
delay, both cited after the proposition. The formal statement takes the family as given, so
these two results are outside the formal development. The examples, Table~\ref{tab:matern},
the list of $c$ in \S\ref{sec:isotropic-origin} and the figures are computed, in the text or by
scripts, and the statements are independent of the computed values. That a kernel with a
Gaussian part is bounded is argued in \S\ref{sec:isotropic-origin} from \ledger{A1}, with no
formal counterpart.
% Q-0450 (the scope audit of 2026-10-13, mismatch 3): the abstract's threshold covers the
% kernels with a Gaussian part, and for them it is the argument of the paragraph before Theorem
% thm:origin-boundedness-d (the transform is at most e^{-a|xi|^2/2}, and C1 applies), which is
% no numbered statement and has no declaration (isotropic_stage_smooth_density is for a stage,
% 0 < s < t). The abstract's last sentence sends the reader here for what is not
% machine-checked, so it is said here.
% trust base: end
